\paragraph{Brain-LLM Alignment} 
A growing body of work has established that the internal representations of large language models (LLMs) exhibit meaningful similarity to neural activity recorded during vision \citep{yamins2014performance, khaligh2014deep, cichy2016comparison, gokce2024scaling}, auditory \citep{koumura2023human, kell2018task, tuckute2023many}, and human language processing \citep{schrimpf2021neural, goldstein2022shared, caucheteux2022brains, tuckute2024driving}. These studies have demonstrated that LLM representations can linearly predict brain responses to the same linguistic stimuli across a variety of neuroimaging paradigms, including reading naturalistic stories and sentences. More recent work has extended these findings by examining how specific model properties relate to brain alignment, finding strong correlations with model size and world knowledge, and demonstrating that instruction-tuning further improves alignment beyond what is achieved by pretraining alone \citep{aw2023instruction}. Parallel efforts have examined how brain alignment evolves over the course of LLM training, revealing that alignment with the human language network saturates early and tracks the acquisition of formal linguistic competence (knowledge of syntactic rules and grammatical structure) more closely than functional competence involving world knowledge and reasoning \citep{alkhamissi2025language}. However, it is notable that prior alignment work has been conducted largely using passive language paradigms in which participants read or listen to naturalistic text. Only recently, some works have examined the brain alignment in active cognitive reasoning tasks such as novel word association generation \citep{kwon2024assessing} and abstract reasoning \citep{pinier2025large}. Closest to our work is \citet{kwon2024assessing}, which has evaluated alignment across two task conditions: A creative condition where participants had to generate original associates to a target word, and a non-creative condition where participants would generate appropriate associations. Findings have revealed that LLM activations align with human brain activity during creative word associations, but not during appropriate word associations, indicating that complex language production is crucial to alignment. However, to the best of our knowledge, no prior work has examined brain alignment during active creative thinking. We address this gap by studying brain–LLM alignment in an active creative task (AUT task), extending emerging efforts that probe alignment in settings requiring cognitive flexibility.

\paragraph{LLMs and Creative Thinking}
Beyond passive language comprehension, LLMs have recently been shown to exhibit impressive performance on tasks requiring creative and divergent thinking \citep{ismayilzada2024creativity}. Studies evaluating models on creativity tasks such as the AUT have found that frontier models can generate responses that are rated as highly original, sometimes matching or exceeding average human performance \citep{goes2023pushing, stevenson2022putting, bellemare2026divergent}. Prior work has proposed various techniques to improve creative thinking and problem-solving in LLMs, such as creative prompting \citep{nguyen2025divergent, tian2024macgyver, lu2025benchmarking} and fine-tuning/preference optimization strategies \citep{puerto2025fine, ismayilzada-etal-2025-creative}. Overall, progress in this field has mainly focused on evaluating and improving divergent creative thinking in LLMs; however, the question of whether high performance on creative thinking tasks is accompanied by representational similarity to the neural systems that support human creativity remains unexplored. We aim to investigate this question of brain-LLM alignment in the present work.